% 第 29 章骨架：导引框已定，正文待写。写作计划见 docs/PLAN.md“第 29 章”，写法见 docs/STYLE.md。
\chapter{大萧条：世界与“金色枷锁”}\label{ch:depressionworld}

\begin{chapterguide}
\guidequestion{金本位怎样把萧条传到全世界？银本位的中国为什么被美国的白银政策拖累？}
\begin{guidegoals}
\item 用“金色枷锁”框架（Eichengreen 1992）解释萧条的国际传导；
\item 比较各国脱离金本位的时间与复苏速度；
\item 分析美国《白银收购法》（1934）对中国的冲击与 1935 年法币改革；
\item 描述拉美与初级产品出口国的经历。
\end{guidegoals}
\guideprereq{\chref{ch:depressionus}\enspace\nameref{ch:depressionus}；\chref{ch:periphery}\enspace\nameref{ch:periphery}}
\end{chapterguide}

\chapterstub
